\chapter{エントロピー正則化}
\label{ch:seminar-03-entropic}

%% ============================================================
%% 3.1 エントロピー正則化
%% ============================================================
\section{エントロピー正則化}
\label{sec:sem-entropic-regularization}

\begin{definition}{離散エントロピー}{sem-discrete-entropy}
  非負行列 $\mathbf{P} \in \R_{\geq 0}^{n \times m}$ に対して，
  \textbf{離散エントロピー}を
  \[
    \Hb(\mathbf{P})
    \defeq
    - \sum_{i=1}^{n} \sum_{j=1}^{m}
      P_{i,j}\bigl(\log P_{i,j} - 1\bigr)
  \]
  で定める．ただし $0\log 0 = 0$ と約束する．
\end{definition}

\begin{remark}{通常の Shannon エントロピーとの違い}{sem-entropy-convention}
  確率行列では $\sum_{i,j}P_{i,j}=1$ なので，
  \[
    \Hb(\mathbf{P})
    =
    -\sum_{i,j} P_{i,j}\log P_{i,j} + 1
  \]
  であり，通常の Shannon エントロピーとは定数 $1$ だけ異なる．
  この定数は最適化の解には影響しない．この形を使う理由は，
  $-\Hb$ の微分が
  \[
    \frac{\partial(-\Hb)}{\partial P_{i,j}} = \log P_{i,j}
  \]
  と簡単になるからである．
\end{remark}

\begin{remark}{$\varphi(x)=x\log x-x$ の連続性}{sem-phi-continuous}
  $-\Hb$ の被加数 $\varphi(x) \defeq x\log x - x$ は $[0,\infty)$ 上で連続である．
  実際，$(0,\infty)$ 上では連続関数の積・和として連続であり，
  $x = 0$ では，約束 $0\log 0 = 0$ により $\varphi(0) = 0$ と定めると，
  $\lim_{x\to0+} x\log x = 0$ から
  \[
    \lim_{x \to 0+} \varphi(x)
    = \lim_{x \to 0+}\bigl(x\log x - x\bigr) = 0 = \varphi(0)
  \]
  となって $x = 0$ でも（右）連続である．
  ゆえに $-\Hb(\mathbf{P}) = \sum_{i,j}\varphi(P_{i,j})$ は
  $\R_{\geq 0}^{n\times m}$ 上で連続である．
\end{remark}

\begin{definition}{エントロピー正則化された離散最適輸送}{sem-entropic-ot}
  $\mathbf{a} \in \R_{>0}^n$（$\sum_{i=1}^n a_i = 1$），
  $\mathbf{b} \in \R_{>0}^m$（$\sum_{j=1}^m b_j = 1$），
  $\mathbf{C} \in \R_{\geq 0}^{n \times m}$，$\varepsilon > 0$ とし，
  \[
    \CouplingsD(\mathbf{a}, \mathbf{b})
    \defeq
    \left\{
      \mathbf{P} \in \R_{\geq 0}^{n \times m}
      \;\middle|\;
      \mathbf{P}\ones_m = \mathbf{a},\;
      \mathbf{P}^\top\ones_n = \mathbf{b}
    \right\}
  \]
  とおく．\textbf{エントロピー正則化された離散 Kantorovich 問題}を
  \[
    \MKD_{\mathbf{C}}^\varepsilon(\mathbf{a}, \mathbf{b})
    \defeq
    \min_{\mathbf{P} \in \CouplingsD(\mathbf{a}, \mathbf{b})}
    \left\{
      \inner{\mathbf{C}}{\mathbf{P}}
      - \varepsilon \Hb(\mathbf{P})
    \right\}
  \]
  と定める．
\end{definition}

\begin{claim}{負エントロピーの狭義凸性}{sem-neg-entropy-strict-convex}
  負エントロピー
  \[
    -\Hb(\mathbf{P}) = \sum_{i,j}\varphi(P_{i,j}),
    \qquad \varphi(x) \defeq x\log x - x
  \]
  は $\CouplingsD(\mathbf{a}, \mathbf{b})$ 上で狭義凸である．
\end{claim}

\begin{proof}
  $-\Hb(\mathbf{P}) = \sum_{i,j}\varphi(P_{i,j})$（$\varphi(x) = x\log x - x$）であるから，
  \textbf{(1)} $\varphi$ が $[0,\infty)$ 上で狭義凸であること，
  \textbf{(2)} そこから和 $\sum_{i,j}\varphi(P_{i,j})$ が狭義凸になること，
  の順に示す．

  \textbf{(1) $\varphi$ の狭義凸性．}
  $\varphi$ は $[0,\infty)$ 上で連続（注意~\ref{rem:sem-phi-continuous}），
  $(0,\infty)$ 上で微分可能で，
  $\varphi'(x) = \log x$ は $(0,\infty)$ 上で狭義単調増加である．
  相異なる $x_1, x_2 \in [0,\infty)$ と $\lambda \in (0,1)$ をとり，
  $x^* \defeq \lambda x_1 + (1-\lambda)x_2 \in (0,\infty)$ とおく．
  目標は
  \[
    \varphi(x^*) < \lambda\varphi(x_1) + (1-\lambda)\varphi(x_2)
  \]
  を示すことである．

  $\varphi$ は閉区間 $[x_1, x^*]$, $[x^*, x_2]$ 上で連続，その内部で微分可能だから，
  平均値定理（定理~\ref{thm:sem-mvt}）より
  \[
    \varphi(x_1) - \varphi(x^*) = \varphi'(c_1)\,(x_1 - x^*),
    \qquad
    \varphi(x_2) - \varphi(x^*) = \varphi'(c_2)\,(x_2 - x^*)
  \]
  をみたす $c_1$（$x_1$ と $x^*$ の間），$c_2$（$x^*$ と $x_2$ の間）が存在する．
  $x_1 = 0$ のときも $c_1 \in (0, x^*)$ は内点なので $\varphi'(c_1) = \log c_1$ は定義される．

  $x^* = \lambda x_1 + (1-\lambda)x_2$ より
  $x_1 - x^* = (1-\lambda)(x_1 - x_2)$,
  $x_2 - x^* = -\lambda(x_1 - x_2)$
  である．これらを上の 2 式に代入し，第 1 式を $\lambda$ 倍，第 2 式を $(1-\lambda)$ 倍して加えると
  \begin{align*}
    \lambda\varphi(x_1) + (1-\lambda)\varphi(x_2) - \varphi(x^*)
    &= \lambda(1-\lambda)(x_1 - x_2)\,\varphi'(c_1)
       - (1-\lambda)\lambda(x_1 - x_2)\,\varphi'(c_2) \\
    &= \lambda(1-\lambda)(x_1 - x_2)\bigl(\varphi'(c_1) - \varphi'(c_2)\bigr)
  \end{align*}
  を得る（左辺は $\lambda + (1-\lambda) = 1$ より $\varphi(x^*)$ がまとまる）．

  最後にこの符号を調べる．$c_1, c_2$ は $x^*$ をはさんで反対側にあるから
  $c_1 - c_2$ は $x_1 - x_2$ と同符号であり，$\varphi'$ が狭義単調増加だから
  $\varphi'(c_1) - \varphi'(c_2)$ も $c_1 - c_2$（ひいては $x_1 - x_2$）と同符号である．
  ゆえに $(x_1 - x_2)\bigl(\varphi'(c_1) - \varphi'(c_2)\bigr) > 0$ となり，
  $\lambda(1-\lambda) > 0$ と合わせて右辺は正．したがって目標の不等式が成り立ち，
  $\varphi$ は $[0,\infty)$ 上で狭義凸である．

  \textbf{(2) $-\Hb$ の狭義凸性．}
  相異なる $\mathbf{P}, \mathbf{Q} \in \CouplingsD(\mathbf{a}, \mathbf{b})$ と $t \in (0,1)$ をとり，
  凸結合 $(1-t)\mathbf{P} + t\mathbf{Q}$ における $-\Hb$ を評価する．
  $-\Hb(\mathbf{R}) = \sum_{i,j}\varphi(R_{i,j})$ であり，
  行列 $(1-t)\mathbf{P} + t\mathbf{Q}$ の $(i,j)$ 成分は $(1-t)P_{i,j} + tQ_{i,j}$ だから，
  \begin{align*}
    -\Hb\bigl((1-t)\mathbf{P} + t\mathbf{Q}\bigr)
    &= \sum_{i,j}\varphi\bigl((1-t)P_{i,j} + tQ_{i,j}\bigr) \\
    &\leq \sum_{i,j}\bigl[(1-t)\varphi(P_{i,j}) + t\varphi(Q_{i,j})\bigr] \\
    &= (1-t)\sum_{i,j}\varphi(P_{i,j}) + t\sum_{i,j}\varphi(Q_{i,j}) \\
    &= (1-t)\bigl(-\Hb(\mathbf{P})\bigr) + t\bigl(-\Hb(\mathbf{Q})\bigr)
  \end{align*}
  となる．等号の各行は $-\Hb = \sum_{i,j}\varphi$ の定義と和の線形性（スカラーのくくり出し）による．
  第 2 行の不等号は，各成分に (1) の狭義凸性
  $\varphi\bigl((1-t)P_{i,j}+tQ_{i,j}\bigr) \leq (1-t)\varphi(P_{i,j}) + t\varphi(Q_{i,j})$
  を適用したもので，$P_{i,j} \neq Q_{i,j}$ の成分では狭義（$<$），
  $P_{i,j} = Q_{i,j}$ の成分では等号である．
  $\mathbf{P} \neq \mathbf{Q}$ より狭義となる成分が少なくとも 1 つあるので，
  上の不等号は全体として狭義であり，
  \[
    -\Hb\bigl((1-t)\mathbf{P} + t\mathbf{Q}\bigr)
    < (1-t)\bigl(-\Hb(\mathbf{P})\bigr) + t\bigl(-\Hb(\mathbf{Q})\bigr).
  \]
  よって $-\Hb$ は $\CouplingsD(\mathbf{a}, \mathbf{b})$ 上で狭義凸である．
\end{proof}

\begin{proposition}{正則化解の存在と一意性}{sem-entropic-existence-unique}
  任意の $\varepsilon > 0$ に対して，
  $\MKD_{\mathbf{C}}^\varepsilon(\mathbf{a}, \mathbf{b})$ は一意な最適解
  $\mathbf{P}_\varepsilon \in \CouplingsD(\mathbf{a}, \mathbf{b})$ を持つ．
\end{proposition}

\begin{proof}
  \textbf{存在．}
  $\CouplingsD(\mathbf{a}, \mathbf{b})$ は空でないコンパクト集合である
  （主張~\ref{clm:sem-discrete-kantorovich-existence} の証明）．
  $\varphi(x) = x\log x - x$ は $[0,\infty)$ 上で連続（注意~\ref{rem:sem-phi-continuous}）だから，
  目的関数
  \[
    F(\mathbf{P})
    \defeq
    \inner{\mathbf{C}}{\mathbf{P}} - \varepsilon \Hb(\mathbf{P})
    =
    \inner{\mathbf{C}}{\mathbf{P}}
    + \varepsilon \sum_{i,j} \varphi(P_{i,j})
  \]
  も連続である．
  コンパクト集合上の連続関数は最小値を達成するから，
  最適解が存在する．

  \textbf{一意性．}
  正則化項がコストの線形性を破り，目的関数 $F$ を狭義凸化する点が鍵である．
  $\mathbf{P}^\star, \mathbf{Q}^\star \in \CouplingsD(\mathbf{a}, \mathbf{b})$ をともに最適解とし，
  最適値を $m \defeq \MKD_{\mathbf{C}}^\varepsilon(\mathbf{a}, \mathbf{b}) = F(\mathbf{P}^\star) = F(\mathbf{Q}^\star)$ とおく．
  $\mathbf{P}^\star \neq \mathbf{Q}^\star$ と仮定する．
  中点 $\mathbf{R} \defeq \tfrac{1}{2}(\mathbf{P}^\star + \mathbf{Q}^\star)$ は，
  $\CouplingsD(\mathbf{a}, \mathbf{b})$ の凸性（主張~\ref{clm:sem-coupling-convex}）よりこれに属する．
  主張~\ref{clm:sem-neg-entropy-strict-convex}（$-\Hb$ の狭義凸性）を
  中点 $\mathbf{R} = \tfrac{1}{2}\mathbf{P}^\star + \tfrac{1}{2}\mathbf{Q}^\star$ に適用すると，
  $\mathbf{P}^\star \neq \mathbf{Q}^\star$ より狭義不等号
  \[
    -\Hb(\mathbf{R})
    < \tfrac{1}{2}\bigl(-\Hb(\mathbf{P}^\star)\bigr) + \tfrac{1}{2}\bigl(-\Hb(\mathbf{Q}^\star)\bigr)
  \]
  が成り立つ．また $\inner{\mathbf{C}}{\cdot}$ は線形だから
  \[
    \inner{\mathbf{C}}{\mathbf{R}}
    = \tfrac{1}{2}\inner{\mathbf{C}}{\mathbf{P}^\star} + \tfrac{1}{2}\inner{\mathbf{C}}{\mathbf{Q}^\star}
  \]
  が等号で成り立つ．これらより（$\varepsilon > 0$ は不等号の向きを保つ）
  \begin{align*}
    F(\mathbf{R})
    &= \inner{\mathbf{C}}{\mathbf{R}} - \varepsilon\Hb(\mathbf{R}) \\
    &= \Bigl(\tfrac{1}{2}\inner{\mathbf{C}}{\mathbf{P}^\star} + \tfrac{1}{2}\inner{\mathbf{C}}{\mathbf{Q}^\star}\Bigr)
       - \varepsilon\Hb(\mathbf{R}) \\
    &< \Bigl(\tfrac{1}{2}\inner{\mathbf{C}}{\mathbf{P}^\star} + \tfrac{1}{2}\inner{\mathbf{C}}{\mathbf{Q}^\star}\Bigr)
       + \tfrac{1}{2}\bigl(-\varepsilon\Hb(\mathbf{P}^\star)\bigr)
       + \tfrac{1}{2}\bigl(-\varepsilon\Hb(\mathbf{Q}^\star)\bigr) \\
    &= \tfrac{1}{2}\bigl(\inner{\mathbf{C}}{\mathbf{P}^\star} - \varepsilon\Hb(\mathbf{P}^\star)\bigr)
       + \tfrac{1}{2}\bigl(\inner{\mathbf{C}}{\mathbf{Q}^\star} - \varepsilon\Hb(\mathbf{Q}^\star)\bigr) \\
    &= \tfrac{1}{2}F(\mathbf{P}^\star) + \tfrac{1}{2}F(\mathbf{Q}^\star) = m
  \end{align*}
  を得る．ここで第 2 行で線形性の等式を $\inner{\mathbf{C}}{\mathbf{R}}$ に代入し，
  第 3 行で $-\varepsilon\Hb(\mathbf{R})$ を狭義凸性（の $\varepsilon$ 倍）で上から押さえた．しかし $\mathbf{R} \in \CouplingsD(\mathbf{a}, \mathbf{b})$ かつ $m$ は $F$ の最小値だから
  $F(\mathbf{R}) \geq m$ でなければならず，矛盾．
  よって $\mathbf{P}^\star = \mathbf{Q}^\star$，すなわち最適解は一意である．
\end{proof}

%% ============================================================
%% 3.2 KL ダイバージェンスによる定式化
%% ============================================================
\section{KL ダイバージェンスによる定式化}
\label{sec:sem-kl-formulation}

\begin{definition}{離散 KL ダイバージェンス}{sem-kl-divergence}
  非負行列 $\mathbf{P}, \mathbf{K} \in \R_{\geq 0}^{n \times m}$ に対して，
  \textbf{KL ダイバージェンス}を
  \[
    \KLD(\mathbf{P} \| \mathbf{K})
    \defeq
    \sum_{i,j}
    \left(
      P_{i,j}\log\frac{P_{i,j}}{K_{i,j}}
      - P_{i,j} + K_{i,j}
    \right)
  \]
  で定める．ただし $0\log 0 = 0$ と約束し，
  $P_{i,j} > 0$ かつ $K_{i,j}=0$ となる成分がある場合は
  $\KLD(\mathbf{P}\|\mathbf{K}) = +\infty$ とする．
\end{definition}

\begin{definition}{Gibbs カーネル}{sem-gibbs-kernel}
  コスト行列 $\mathbf{C}$ と $\varepsilon > 0$ に対して，
  \textbf{Gibbs カーネル} $\mathbf{K} \in \R_{>0}^{n \times m}$ を
  \[
    K_{i,j} \defeq \exp\!\left(-\frac{C_{i,j}}{\varepsilon}\right)
  \]
  で定める．
\end{definition}

\begin{proposition}{正則化 OT は KL 射影である}{sem-entropic-kl-projection}
  $\mathbf{K} = \exp(-\mathbf{C}/\varepsilon)$ とすると，
  $\MKD_{\mathbf{C}}^\varepsilon(\mathbf{a}, \mathbf{b})$ の最適解は
  \[
    \mathbf{P}_\varepsilon
    =
    \argmin_{\mathbf{P} \in \CouplingsD(\mathbf{a}, \mathbf{b})}
    \KLD(\mathbf{P}\|\mathbf{K})
  \]
  と書ける．
\end{proposition}

\begin{proof}
  $\log K_{i,j} = -C_{i,j}/\varepsilon$ より
  \begin{align*}
    \varepsilon \KLD(\mathbf{P}\|\mathbf{K})
    &=
    \varepsilon\sum_{i,j}
    \left(
      P_{i,j}\log P_{i,j}
      - P_{i,j}\log K_{i,j}
      - P_{i,j}
      + K_{i,j}
    \right) \\
    &=
    \sum_{i,j} C_{i,j}P_{i,j}
    + \varepsilon\sum_{i,j} P_{i,j}(\log P_{i,j}-1)
    + \varepsilon\sum_{i,j} K_{i,j} \\
    &=
    \inner{\mathbf{C}}{\mathbf{P}} - \varepsilon \Hb(\mathbf{P})
    + \varepsilon\sum_{i,j} K_{i,j}.
  \end{align*}
  最後の項は $\mathbf{P}$ に依存しない定数なので，
  両者の最小化問題は同じ最適解を持つ．
\end{proof}

%% ============================================================
%% 3.3 正則化パラメータの極限
%% ============================================================
\section{正則化パラメータの極限}
\label{sec:sem-epsilon-limits}

非正則化 Kantorovich 問題が複数の最適解を持つとき，
エントロピー正則化はそのうちの一つ——
エントロピーが最大のもの——を自然に選び出す．

\begin{theorem}{$\varepsilon \to 0$ による非正則化 OT への収束}{sem-entropic-convergence-eps-zero}
  各 $\varepsilon > 0$ に対し，$\mathbf{P}_\varepsilon$ を
  $\MKD_{\mathbf{C}}^\varepsilon(\mathbf{a}, \mathbf{b})$ の一意解とする．
  $\varepsilon \to 0$ のとき，以下が成り立つ．

  \begin{enumerate}
    \item[\textrm{(i)}]
      最適値は収束する：
      $\displaystyle
        \lim_{\varepsilon \to 0}
        \MKD_{\mathbf{C}}^\varepsilon(\mathbf{a},\mathbf{b})
        =
        \MKD_{\mathbf{C}}(\mathbf{a},\mathbf{b}).
      $
    \item[\textrm{(ii)}]
      $\mathbf{P}_\varepsilon$ の $\varepsilon \to 0$ における任意の極限点は
      非正則化問題の最適解である．
    \item[\textrm{(iii)}]
      $\{\mathbf{P}_\varepsilon\}_{\varepsilon > 0}$ 全体が，
      非正則化最適解の中でエントロピー $\Hb$ を最大化する
      一意な解に収束する：
      \[
        \mathbf{P}_\varepsilon
        \;\xrightarrow{\varepsilon \to 0}\;
        \argmax_{\mathbf{P} \in S^*} \Hb(\mathbf{P}).
      \]
  \end{enumerate}
\end{theorem}

\begin{proof}
  非正則化問題の任意の最適解を $\mathbf{P}^*$ とする．

  \textbf{(i), (ii).}
  $\mathbf{P}_\varepsilon$ の最適性より
  \[
    \inner{\mathbf{C}}{\mathbf{P}_\varepsilon}
    - \varepsilon\Hb(\mathbf{P}_\varepsilon)
    \leq
    \inner{\mathbf{C}}{\mathbf{P}^*}
    - \varepsilon\Hb(\mathbf{P}^*).
  \]
  一方，$\mathbf{P}^*$ は非正則化問題の最適解なので
  $\inner{\mathbf{C}}{\mathbf{P}^*}
  \leq \inner{\mathbf{C}}{\mathbf{P}_\varepsilon}$ であり，
  \[
    0
    \leq
    \inner{\mathbf{C}}{\mathbf{P}_\varepsilon}
    - \inner{\mathbf{C}}{\mathbf{P}^*}
    \leq
    \varepsilon\bigl(\Hb(\mathbf{P}_\varepsilon)-\Hb(\mathbf{P}^*)\bigr).
  \]
  $\CouplingsD(\mathbf{a}, \mathbf{b})$ はコンパクトであり，$\Hb$ はその上で有界なので，
  右辺は $\varepsilon \to 0$ で $0$ に収束する．
  したがって
  $\inner{\mathbf{C}}{\mathbf{P}_\varepsilon}
  \to \inner{\mathbf{C}}{\mathbf{P}^*}
  = \MKD_{\mathbf{C}}(\mathbf{a},\mathbf{b})$ であり，
  $\varepsilon\Hb(\mathbf{P}_\varepsilon) \to 0$ なので
  $\MKD_{\mathbf{C}}^\varepsilon \to \MKD_{\mathbf{C}}$ を得る．
  コンパクト集合 $\CouplingsD(\mathbf{a},\mathbf{b})$ 内の
  任意の極限点は非正則化問題の最適値を達成するから，
  最適解である．

  \textbf{(iii).}
  上の不等式から
  $\Hb(\mathbf{P}_\varepsilon) \geq \Hb(\mathbf{P}^*)$
  が任意の非正則化最適解 $\mathbf{P}^*$ について成り立つ．
  $\Hb$ は連続であるから，任意の極限点 $\mathbf{P}^0$ は
  非正則化最適解の中で $\Hb$ を最大化する．
  最適解集合 $S^*$ は凸であり（主張~\ref{clm:sem-optimal-face}），
  主張~\ref{clm:sem-neg-entropy-strict-convex} より $-\Hb$ は
  （凸部分集合 $S^*$ 上でも）狭義凸であるから，
  命題~\ref{prop:sem-strict-convex-unique-min} より
  $\Hb$ を最大化する $S^*$ の元は一意である．
  したがって $\{\mathbf{P}_\varepsilon\}$ の任意の収束部分列は
  同一の極限を持ち，コンパクト集合内の点列でこの性質を持つものは
  全体として収束するから，
  $\mathbf{P}_\varepsilon$ は最大エントロピー最適解に収束する．
\end{proof}

\begin{remark}{$\varepsilon \to +\infty$ による独立カップリングへの収束}{sem-entropic-convergence-eps-infty}
  $\varepsilon \to +\infty$ のとき
  $\mathbf{P}_\varepsilon \to \mathbf{a}\mathbf{b}^\top$ が成り立つ．
  $\varepsilon$ で目的関数を割ると
  $\varepsilon^{-1}\inner{\mathbf{C}}{\mathbf{P}} - \Hb(\mathbf{P})$
  の最小化であり，$\varepsilon \to +\infty$ では
  エントロピー最大化
  $\max_{\mathbf{P} \in \CouplingsD} \Hb(\mathbf{P})$
  に帰着する．周辺分布を固定したエントロピー最大化の解は
  独立カップリング $\mathbf{a}\mathbf{b}^\top$ である．
\end{remark}
